本问的优点不在于继续增加模型维度，而在于把流域平均单区模型收紧成“结构参数设定--单硬锚点校准--局地一致性检验--反事实对照”的实证闭环。具体而言，2025 年四大家鱼恢复到禁渔前约 1.8 倍是唯一参与拟合的硬校准锚点；\(consumer\_scale=0.6\) 作为归一化初值设定固定给定，不再与 2025 年锚点共同搜索；湖北段单网次捕获量提升 120\% 对应的 \(\mathrm{CPUE}^*\) 仅作为拟合外的局地软约束。按这一口径，校准后四大家鱼总量倍率达到 1.805、相对误差为 0.26\%，而 \(\mathrm{CPUE}^*\) 达到 1.838、相对误差为 16.47\%，说明模型在全流域总量层面已较好贴合题面事实，但对局地捕捞代理只能做到方向一致。与此同时，\(H_0=0.25\) 仅作为区间 \([0.10,0.40]\) 内的固定中位反事实参数，不参与任何观测拟合；在该对照下，2025 年四大家鱼总量和 \(\mathrm{CPUE}^*\) 分别为 0.716 和 0.731，且轨迹无反常上冲，因此“禁渔优于不禁渔”不再依赖目标导向的对照参数搜索。

\(\Phi(t)\) 与水草末值的作用也应收紧理解：它们不是外部观测，而是用于解释“恢复与承压并行”的内部机制指标。当前参数组下，食压指数升至初值的 4.213 倍，水草末值降至 0.204 倍；对 \(gain\_scale\)、\(consumer\_scale\) 和 \(H_0\) 做 \(\pm10\%\) 扰动后，禁渔下鱼群恢复强于不禁渔、且系统承压上升这一主结论方向不变。其局限在于 \(\mathrm{CPUE}^*\) 的 \((1,1,1,0.7)\) 权重属于经验假设，未单独估计，单区均质假设也无法直接解释湖北段等局地空间差异，因此本文不将湖北段观测写成全流域硬拟合结果。后续补强应优先增加独立观测锚点、统一代理口径，并继续补做更广范围的敏感性检验；空间异质性建模只能作为后续扩展，而不是当前版本的直接结论。
